\section*{Chapitre \ref{ch:g4:clouds-rain-snow} --- Les nuages, la pluie et la neige}

\begin{solution}{exo:g4:clouds-rain-snow:1}
Un \omterm{def:g4:clouds-rain-snow:cloud}{nuage} est fait d'un très grand nombre de minuscules gouttelettes d'eau,
ou de minuscules morceaux de glace, qui flottent dans l'air. Le \omterm{def:g4:clouds-rain-snow:cloud}{brouillard}
est un \omterm{def:g4:clouds-rain-snow:cloud}{nuage} qui touche le sol.
\end{solution}

\begin{solution}{exo:g4:clouds-rain-snow:2}
(a) Un \omterm{def:g4:clouds-rain-snow:cloud-families}{stratus}. (b) Des \omterm{def:g4:clouds-rain-snow:cloud-families}{cirrus}. (c) Des \omterm{def:g4:clouds-rain-snow:cloud-families}{cumulus}. (d) Un \omterm{def:g4:clouds-rain-snow:cloud-families}{cumulonimbus}.
\end{solution}

\begin{solution}{exo:g4:clouds-rain-snow:3}
(a) De la \omterm{def:g4:clouds-rain-snow:precipitation}{grêle}. (b) De la \omterm{def:g4:clouds-rain-snow:precipitation}{pluie}. (c) De la \omterm{def:g4:clouds-rain-snow:precipitation}{neige}.
\end{solution}

\begin{solution}{exo:g4:clouds-rain-snow:4}
(a) Un \omterm{def:g4:clouds-rain-snow:rain-gauge}{pluviomètre}. (b) Un thermomètre. (c) Une \omterm{def:g4:clouds-rain-snow:wind-vane}{girouette}.
\end{solution}

\begin{solution}{exo:g4:clouds-rain-snow:5}
17 millimètres de \omterm{def:g4:clouds-rain-snow:precipitation}{pluie} sont tombés. $20 - 17 = 3$~: 3 millimètres de plus
amèneraient l'eau à la graduation 20.
\end{solution}

\begin{solution}{exo:g4:clouds-rain-snow:6}
Le jour le plus pluvieux a été mercredi (12 millimètres). Il n'a plu ni
lundi ni vendredi. Jeudi a eu 7 millimètres.
\end{solution}

\begin{solution}{exo:g4:clouds-rain-snow:7}
Les flocons se forment en haut, dans le \omterm{def:g4:clouds-rain-snow:cloud}{nuage} froid. Là-haut, au village de
\omterm{def:g1:landscapes-around-us:landform}{montagne}, l'air est sous zéro jusqu'en bas, donc ils arrivent sous forme de
\omterm{def:g4:clouds-rain-snow:precipitation}{neige}. Plus bas, dans la \omterm{def:g1:landscapes-around-us:landform}{vallée}, l'air est plus chaud que \SI{0}{\celsius}~:
en le traversant pendant leur chute, les flocons fondent, et ils arrivent
au sol sous forme de \omterm{def:g4:clouds-rain-snow:precipitation}{pluie}.
\end{solution}

\begin{solution}{exo:g4:clouds-rain-snow:8}
Le diagramme a sept bâtons de hauteurs 3, 0, 1, 9, 6, 0 et 1 (pas de bâton
mardi ni samedi). Dans la semaine, $3 + 0 + 1 + 9 + 6 + 0 + 1 = 20$
millimètres sont tombés. Jeudi a apporté 9 des 20 millimètres~:
$\frac{9}{20}$ de la \omterm{def:g4:clouds-rain-snow:precipitation}{pluie} de la semaine.
\end{solution}

\begin{solution}{exo:g4:clouds-rain-snow:9}
Les murs et les arbres abritent de la \omterm{def:g4:clouds-rain-snow:precipitation}{pluie} le sol qui les entoure, surtout
quand le vent pousse la \omterm{def:g4:clouds-rain-snow:precipitation}{pluie} de côté, et des gouttes peuvent tomber des
feuilles dans le \omterm{def:g4:clouds-rain-snow:rain-gauge}{pluviomètre}~: près d'eux, le \omterm{def:g4:clouds-rain-snow:rain-gauge}{pluviomètre} recueillerait
trop peu d'eau, ou trop. Le thermomètre doit mesurer la température de
l'air~; en plein soleil, il serait chauffé par le soleil lui-même et
indiquerait une valeur trop haute. L'abri blanc à lamelles arrête le
soleil et laisse passer l'air.
\end{solution}

\begin{solution}{exo:g4:clouds-rain-snow:10}
Une \omterm{def:g4:clouds-rain-snow:wind-vane}{girouette} montre d'où vient le vent~: le vent vient de l'ouest. Il va
vers l'est.
\end{solution}

\begin{solution}{exo:g4:clouds-rain-snow:11}
Un millimètre de \omterm{def:g4:clouds-rain-snow:precipitation}{pluie} sur un mètre carré, c'est un litre. Sur un mètre
carré, 8 millimètres donnent 8 litres. Sur 500 mètres carrés~:
$500 \times 8 = 4000$ litres. Avec des arrosoirs de 10 litres~:
$4000 \div 10 = 400$ arrosoirs.
\end{solution}
